\documentclass[11pt,a4paper]{ctexart}
\usepackage[a4paper,margin=1.78cm,headheight=15pt]{geometry}
\usepackage{amsmath,amssymb,mathtools,bm}
\usepackage{booktabs,longtable,tabularx,array}
\usepackage{enumitem}
\usepackage[most]{tcolorbox}
\usepackage{tikz}
\usetikzlibrary{arrows.meta,positioning,fit,calc}
\usepackage{xcolor}
\usepackage{fancyhdr}
\usepackage{hyperref}
\usepackage{microtype}

\newcolumntype{Y}{>{\raggedright\arraybackslash}X}
\newcolumntype{L}[1]{>{\raggedright\arraybackslash}p{#1}}
\setlength{\parindent}{2em}
\setlength{\parskip}{0.34em}
\setlength{\emergencystretch}{3em}
\renewcommand{\arraystretch}{1.22}
\setlength{\tabcolsep}{3pt}
\setlist[itemize]{itemsep=0.18em,topsep=0.35em,leftmargin=2em}
\setlist[enumerate]{itemsep=0.18em,topsep=0.35em,leftmargin=2em}

\definecolor{deep}{RGB}{42,58,73}
\definecolor{soft}{RGB}{245,247,249}
\definecolor{line}{RGB}{145,154,163}
\definecolor{accent}{RGB}{104,76,55}
\definecolor{greenish}{RGB}{57,92,76}
\definecolor{warn}{RGB}{136,68,63}
\definecolor{purple}{RGB}{87,72,116}

\hypersetup{
  colorlinks=true,
  linkcolor=deep,
  urlcolor=accent,
  pdftitle={秩、基本子空间与等价}
}

\pagestyle{fancy}
\fancyhf{}
\lhead{\small\color{deep}\textbf{数学一 · 线性代数 Topic 03}}
\rhead{\small\color{deep}秩 · Kernel / Image · 等价}
\cfoot{\small\thepage}
\renewcommand{\headrulewidth}{0.4pt}

\newtcolorbox{corebox}[1]{breakable,colback=soft,colframe=deep,title=\textbf{#1},boxrule=0.8pt,arc=2mm}
\newtcolorbox{methodbox}[1]{breakable,colback=white,colframe=greenish,title=\textbf{#1},boxrule=0.7pt,arc=1.5mm}
\newtcolorbox{warnbox}[1]{breakable,colback=white,colframe=warn,title=\textbf{#1},boxrule=0.7pt,arc=1.5mm}
\newtcolorbox{examplebox}[1]{breakable,colback=white,colframe=purple,title=\textbf{#1},boxrule=0.7pt,arc=1.5mm}
\newcommand{\kw}[1]{\textbf{\color{deep}#1}}

\tikzset{
  nodebox/.style={draw=deep,rounded corners=2mm,text width=3.0cm,minimum height=1.0cm,align=center,fill=soft,font=\small,inner sep=4pt},
  flow/.style={-{Latex[length=2mm]},thick,draw=deep},
  dashedflow/.style={-{Latex[length=2mm]},thick,dashed,draw=purple}
}

\title{\vspace{-1.1cm}\textbf{秩、基本子空间与等价}\\[0.4em]
\large 一份输入自由度，怎样分成“留下的”和“被抹掉的”}
\author{数学一 · 线性代数 Topic 03}
\date{}

\begin{document}
\maketitle

\begin{corebox}{母问题：一个线性作用到底保留了多少独立自由度，又把多少自由度压掉了？}
给一个矩阵 $A$，只说“它有多少个非零元素”几乎没有结构意义。真正重要的是：输入空间中有多少彼此独立的方向能够在输出端留下可区分的信息，又有多少方向被映射完全抹掉。

因此本册的母模型不是“rank 有几种算法”，而是
\[
\boxed{
\text{输入自由度}
=\text{被抹除的自由度}
+\text{留下可见输出的自由度}
}
\]
在线性映射 $T:V\to W$ 上，这句话精确写成
\[
\boxed{
\dim V=\dim\ker T+\dim\operatorname{Im}T.
}
\]
其中
\[
\operatorname{rank}T=\dim\operatorname{Im}T,
\qquad
\operatorname{nullity}T=\dim\ker T.
\]
\end{corebox}

\section{本册负责什么}

\begin{itemize}
  \item \kw{Owns：}rank、kernel/image 的自由度机制、rank--nullity、行秩=列秩、主元/最大非零子式等 rank 观测方式、矩阵等价、rank 运算规律、乘积的中间瓶颈、低秩结构、伴随矩阵的 rank 分层。
  \item \kw{Uses：}Topic 01 的 span/基/维数/正交；Topic 02 的线性映射、两端换基、初等矩阵、determinant 与 adjugate 恒等式。
  \item \kw{Does not own：}给定 $b$ 的相容性与完整解集；相似和特征结构；二次型惯性。
\end{itemize}

Topic 04 会直接调用本册提供的两个输出：
\[
\boxed{b\in\operatorname{Im}A?}
\qquad
\boxed{\dim\ker A=n-r(A).}
\]
但“某个具体 $b$ 是否可达，以及所有逆像怎样写”属于 Topic 04。

\section{先认两个主角：Kernel 与 Image}

设
\[
A\in\mathbb R^{m\times n},
\qquad
A:\mathbb R^n\to\mathbb R^m.
\]

\kw{Kernel / 零空间}是
\[
\ker A=\{x\in\mathbb R^n:Ax=0\}.
\]
它记录所有被作用完全压成零的输入方向。

\kw{Image / 像空间}是
\[
\operatorname{Im}A=\{Ax:x\in\mathbb R^n\}.
\]
若把 $A$ 的列写成 $a_1,\dots,a_n$，则
\[
Ax=x_1a_1+\cdots+x_na_n,
\]
所以
\[
\boxed{\operatorname{Im}A=\operatorname{Col}(A).}
\]

于是矩阵的秩可以直接定义为
\[
\boxed{r(A)=\dim\operatorname{Im}A=\dim\operatorname{Col}(A).}
\]
它不是“非零元素个数”，而是输出端真正能够留下的独立方向数。

\begin{examplebox}{母例 1：二维输入被压成一维}
取
\[
A=\begin{pmatrix}1&1\\1&1\end{pmatrix}.
\]
所有输出都形如
\[
Ax=(x_1+x_2)\begin{pmatrix}1\\1\end{pmatrix},
\]
所以
\[
\operatorname{Im}A
=\operatorname{span}\left\{\begin{pmatrix}1\\1\end{pmatrix}\right\},
\qquad r(A)=1.
\]
同时
\[
Ax=0\iff x_1+x_2=0,
\]
故
\[
\ker A
=\operatorname{span}\left\{\begin{pmatrix}1\\-1\end{pmatrix}\right\}.
\]
二维输入里，一个自由方向留下可见输出，一个自由方向被压掉：
\[
\boxed{2=1+1.}
\]
\end{examplebox}

\section{Rank--Nullity 为什么必然成立}

不要把
\[
n=r(A)+\dim\ker A
\]
当成一个需要背的维数公式。它可以直接从“给 kernel 选基，再把它扩充成整个输入空间的基”推出。

设
\[
k_1,\dots,k_s
\]
是 $\ker T$ 的一组基。把它扩充为 $V$ 的一组基：
\[
k_1,\dots,k_s,v_1,\dots,v_r.
\]
于是
\[
\dim V=s+r.
\]

现在看 $T(v_1),\dots,T(v_r)$。

\subsection{它们生成整个 Image}
任意 $x\in V$ 可写成
\[
x=\sum_{j=1}^{s}b_jk_j+\sum_{i=1}^{r}a_iv_i.
\]
作用 $T$ 后，kernel 部分全部消失：
\[
T(x)=\sum_{i=1}^{r}a_iT(v_i).
\]
所以 $T(v_1),\dots,T(v_r)$ 生成 $\operatorname{Im}T$。

\subsection{它们还线性无关}
若
\[
\sum_{i=1}^{r}a_iT(v_i)=0,
\]
则
\[
T\left(\sum_{i=1}^{r}a_iv_i\right)=0,
\]
所以 $\sum a_iv_i\in\ker T$。但
\[
k_1,\dots,k_s,v_1,\dots,v_r
\]
是一组基，表示唯一；一个只含 $v_i$ 的组合若又落在由 $k_j$ 张成的 kernel 中，只能所有 $a_i=0$。

因此 $T(v_1),\dots,T(v_r)$ 是 $\operatorname{Im}T$ 的一组基，故
\[
\dim\operatorname{Im}T=r.
\]
最终得到
\[
\boxed{
\dim V
=\dim\ker T+\dim\operatorname{Im}T.
}
\]

\begin{corebox}{Rank--Nullity 的真正含义}
它不是说每根输入基向量单独“选择留下或消失”。真正的结构是：\kw{可以选择一组适合 kernel 的输入基，使一部分基方向正好张成 kernel，其余方向的像恰好成为 image 的基。}

因此“被抹除自由度 + 可见自由度 = 输入自由度”是一份可以通过选基显式看见的分解。
\end{corebox}

\section{同一个 rank 为什么会有这么多观测方式}

对 $A\in\mathbb R^{m\times n}$，下面这些数字最终都相同：

\begin{longtable}{L{4.2cm}L{\dimexpr\linewidth-4.7cm\relax}}
\toprule
\textbf{观测方式} & \textbf{它在看什么}\\
\midrule
列向量组的极大无关组个数 & 输出端能生成多少独立方向\\
$\dim\operatorname{Col}(A)$ & image 的维数\\
$\dim\operatorname{Row}(A)$ & 独立行约束 / 输入线性函数有多少独立方向\\
阶梯形中的主元个数 & 消元后剩下多少独立结构\\
最大非零子式的阶数 & 最大能截出多大的可逆方块\\
线性映射的 rank & 作用后保留下来的自由度\\
\bottomrule
\end{longtable}

这些不是六个互相独立的定义，而是同一份“独立信息维数”在向量、空间、消元、determinant 和映射语言中的不同投影。

\section{为什么行秩等于列秩}

对矩阵做可逆行变换。Topic 02 已说明：左乘可逆矩阵不会改变列之间的线性关系，因此列 rank 不变；同时新旧行互为可逆线性组合，所以行空间维数也不变。

把 $A$ 化为行阶梯形 $R$。若 $R$ 有 $r$ 个非零行，则：

\begin{itemize}
  \item 非零行的首非零位置逐级右移，因此这 $r$ 行线性无关，行 rank 为 $r$；
  \item 恰有 $r$ 个主元列，主元列线性无关，其余列由它们生成，因此列 rank 也为 $r$。
\end{itemize}

于是
\[
\boxed{\dim\operatorname{Row}(A)=\dim\operatorname{Col}(A)=r(A).}
\]

\begin{warnbox}{行空间和列空间不是同一个集合}
它们只有\kw{维数相等}。若 $A$ 是 $m\times n$ 非方阵，
\[
\operatorname{Row}(A)\subseteq\mathbb R^n,
\qquad
\operatorname{Col}(A)\subseteq\mathbb R^m,
\]
甚至生活在不同环境空间里，不能写成“行空间=列空间”。
\end{warnbox}

\section{最大非零子式与主元：为什么它们也能测 rank}

一个 $r$ 阶子式非零，表示从 $A$ 中选出的 $r$ 行、$r$ 列组成一个可逆方阵，因此至少存在 $r$ 个独立方向：
\[
r(A)\ge r.
\]

如果所有 $(r+1)$ 阶子式都为零，就不可能再截出一个 $(r+1)$ 维可逆作用，因此
\[
r(A)\le r.
\]
所以
\[
\boxed{
r(A)=r
\iff
\begin{cases}
\text{至少一个 $r$ 阶子式非零},\\
\text{所有 $(r+1)$ 阶子式为零}.
\end{cases}}
\]

这也是 determinant 与 rank 的责任分界：对 $n\times n$ 方阵，$\det A\neq0$ 只说明 $r(A)=n$；若 $\det A=0$，具体是 $n-1$ 还是更低，必须继续由 rank 判断。

\section{矩阵等价：把 kernel 与 image 都换到最干净的坐标}

Rank--Nullity 的基扩充证明还能继续生成一件更强的事。

对 $T:V_n\to W_m$，沿用 Rank--Nullity 中得到的两组方向，但为了让标准形的单位块出现在左上角，把输入基\kw{重新排序}为
\[
v_1,\dots,v_r,k_1,\dots,k_{n-r},
\]
其中 $k_j$ 张成 kernel，而 $T(v_1),\dots,T(v_r)$ 是 image 的一组基。再在输出端把
\[
T(v_1),\dots,T(v_r)
\]
扩充成 $W$ 的一组基。

在这两组“适配映射结构”的基下，前 $r$ 个输入基向量依次送到前 $r$ 个输出基向量，后 $n-r$ 个 kernel 基向量全部送到零，因此 $T$ 的矩阵直接变成
\[
\boxed{
J_r=
\begin{pmatrix}
I_r&0\\
0&0
\end{pmatrix}_{m\times n}.
}
\]

这张矩阵把映射的全部本质压缩成一句话：\kw{前 $r$ 个输入方向被无损保留，其余 $n-r$ 个输入方向被压成零。}

一般矩阵等价写成
\[
B=PAQ,
\]
其中 $P,Q$ 可逆。它允许输入端和输出端独立重编码，因此
\[
\boxed{A\sim_{\mathrm{eq}}B\iff r(A)=r(B)}
\]
（同型矩阵）。

原因已经不是“初等变换碰巧保持 rank”，而是：每个 rank 为 $r$ 的映射都能选基化成同一个 $J_r$。

\begin{warnbox}{同 rank 只保证等价，不保证相似}
矩阵等价允许两端独立换基；相似要求同一空间算子在输入/输出端同步换同一组基。因此
\[
\text{同 rank}\Longrightarrow\text{等价},
\]
但一般不能推出相似。

例如 $I_2$ 与 $\operatorname{diag}(1,2)$ 都 rank 2，因此等价；但特征值不同，所以不相似。
\end{warnbox}

\section{初等变换到底保护哪些子空间}

设 $P,Q$ 可逆。

\begin{longtable}{L{2.7cm}L{4.1cm}L{\dimexpr\linewidth-7.3cm\relax}}
\toprule
\textbf{变换} & \textbf{稳定的东西} & \textbf{会怎样移动}\\
\midrule
$A\mapsto PA$ & rank、列间线性关系、row space & $\operatorname{Col}(PA)=P\operatorname{Col}(A)$，具体输出方向会移动\\
$A\mapsto AQ$ & rank、column space & $\ker(AQ)=Q^{-1}\ker A$，row space 在输入坐标中被 $Q$ 重编码\\
$A\mapsto PAQ$ & rank & 输入/输出两侧的具体方向都可以变，只保留自由度数量\\
\bottomrule
\end{longtable}

因此“初等变换保持什么”必须带对象：\kw{rank 始终保持；具体子空间的位置未必保持。}

\subsection{分块 rank：把普通初等变换提升到块层面}

分块矩阵里出现单位块或可逆主块时，不要把“大块”重新展开成逐元素消元。块行 / 块列变换仍然只是左乘 / 右乘可逆矩阵，因此同样保持 rank。

最常用的可逆块初等矩阵是
\[
L=\begin{pmatrix}I&0\\-Y&I\end{pmatrix},
\qquad
R=\begin{pmatrix}I&-X\\0&I\end{pmatrix},
\]
它们的逆只需把 $-Y,-X$ 改成 $Y,X$。于是对尺寸相容的块矩阵
\[
M=\begin{pmatrix}I&X\\Y&Z\end{pmatrix},
\]
先做块行消元，再做块列消元：
\[
LMR
=
\begin{pmatrix}
I&0\\
0&Z-YX
\end{pmatrix}.
\]
因此
\[
\boxed{
r(M)=r(I)+r(Z-YX).
}
\]

这里的核心不是记一条新公式，而是识别：\kw{单位块给出了一个不会丢失自由度的主元，可以用它把同一行 / 列上的其他块消掉。}

更一般地，若左上块 $A$ 可逆，同样的块消元给出
\[
\boxed{
r\begin{pmatrix}A&B\\C&D\end{pmatrix}
=r(A)+r(D-CA^{-1}B).}
\]
右端的 $D-CA^{-1}B$ 是消去 $A$ 所负责方向以后剩余的有效作用。这里仅把它作为 rank 消元接口；一般 Schur 补理论不进入数学一主干。

\begin{warnbox}{块消元也必须满足乘法顺序与尺寸}
块不是标量。消去左下块时出现的是 $CA^{-1}$，消去右上块时出现的是 $A^{-1}B$；不能交换次序。若主块不可逆，也不能为了套公式写 $A^{-1}$，应改用其他可逆块、普通初等变换或已有零块结构。
\end{warnbox}

\section{四个基本子空间：主干只有一对，另外两个是正交镜像}

在标准内积下，矩阵还自然带出 row space 与 $N(A^T)$：
\[
\operatorname{Row}(A)=\operatorname{Col}(A^T).
\]

对任意 $x$，$Ax=0$ 等价于 $x$ 与 $A$ 的每一行都正交，因此
\[
\boxed{\ker A=\operatorname{Row}(A)^\perp.}
\]
同理
\[
\boxed{\ker A^T=\operatorname{Col}(A)^\perp.}
\]
于是
\[
\mathbb R^n=\operatorname{Row}(A)\oplus\ker A,
\]
\[
\mathbb R^m=\operatorname{Col}(A)\oplus\ker A^T.
\]

\begin{methodbox}{数学一里怎样给这四个空间排优先级？}
真正反复生成题目的是
\[
\boxed{\ker A\quad\text{和}\quad\operatorname{Im}A=\operatorname{Col}(A).}
\]
row space 主要用于解释行秩、独立约束和齐次同解；$\ker A^T$ 主要解释“输出空间中与所有可达方向正交的部分”。

所以不要把“四大基本子空间”再学成四份并列名词表。它们是 kernel/image 主模型在标准内积下的两组正交镜像。
\end{methodbox}

\section{Rank-one：最小的非零信息通道}

令
\[
A=uv^T,
\qquad u\neq0,\ v\neq0.
\]
则
\[
Ax=u(v^Tx).
\]
输入 $x$ 先被压缩成一个标量 $v^Tx$，再沿 $u$ 方向输出。因此
\[
\operatorname{Im}A=\operatorname{span}\{u\},
\qquad r(A)=1,
\]
而
\[
\ker A=\{x:v^Tx=0\}=v^\perp.
\]

这给“低秩”的最小直觉：\kw{低秩不是元素少，而是作用只通过少量独立标量通道传递信息。}

若进一步是方阵 $A=uv^T\in\mathbb R^{n\times n}$，记
\[
g=v^Tu=\operatorname{tr}A.
\]
那么外积可以直接闭合：
\[
A^2=uv^Tuv^T=(v^Tu)uv^T=gA,
\]
因此
\[
A^1=A,
\qquad
\boxed{A^k=g^{k-1}A\qquad(k\ge2).}
\]
这样写可以避免 $g=0$、$k=1$ 时出现 $0^0$ 的形式歧义；真正的结构结论仍是：从二次幂开始，所有高次幂都压回同一条 rank-one 通道。

\begin{corebox}{方阵 rank-one 的迹为什么是谱接口？}
因为 $Au=gu$，所以 $g$ 是一个特征值；而
\[
\ker A=v^\perp
\]
维数为 $n-1$，因此 $0$ 至少提供 $n-1$ 个独立特征方向。

于是：
\begin{itemize}
  \item 若 $g\neq0$，$\mathbb R^n=\operatorname{span}\{u\}\oplus\ker A$，$A$ 可对角化，谱为 $g,0,\dots,0$；
  \item 若 $g=0$ 且 $A\neq0$，则 $A^2=0$，全部特征值都是 $0$，但 $\ker A$ 只有 $n-1$ 维，所以不可对角化。
\end{itemize}
完整“多项式关系为什么控制可对角化”的机制由 Topic 05 接管；这里把 rank-one 作为跨册母例输出。
\end{corebox}

\begin{methodbox}{怎样从数表识别 rank one？}
对非零矩阵，下列信号都在表达同一件事：所有输出只剩一条独立方向。
\begin{itemize}
  \item 所有非零列彼此成比例；
  \item 所有非零行彼此成比例；
  \item 可以直接写成外积 $uv^T$；
  \item 所有 $2\times2$ 子式都为 $0$。
\end{itemize}
最后一条只能先推出 $r(A)\le1$，还必须排除 $A=0$ 才能断言 $r(A)=1$。
\end{methodbox}

\begin{examplebox}{Extension：奇异值分解 (SVD) 的高级镜头}
Rank-one 矩阵的外积结构 $A = uv^T$ 实际上是所有矩阵最底层的“乐高积木”。根据奇异值分解 (SVD) 定理，任何矩阵都可以写成一系列 rank-one 矩阵的加权和：
\[
A = \sigma_1 u_1 v_1^T + \sigma_2 u_2 v_2^T + \cdots + \sigma_r u_r v_r^T.
\]
这里 $\sigma_i$ 是非零奇异值，$u_i$ 和 $v_i$ 分别是输出空间和输入空间的一组标准正交基。
这给出了 rank 的终极统一视角：\kw{矩阵的 rank 恰好等于其非零奇异值的个数}，而每一个 rank-one 分量都代表着一个独立的信息传输通道。
\end{examplebox}

\section{乘积的 rank：真正的瓶颈是 $\operatorname{Im}B$ 撞上 $\ker A$}

设
\[
B:\mathbb R^p\to\mathbb R^n,
\qquad
A:\mathbb R^n\to\mathbb R^m.
\]
复合为 $AB$。先由 $B$ 把输入送进 $\operatorname{Im}B$，再由 $A$ 作用。真正会在第二段继续消失的，是
\[
\operatorname{Im}B\cap\ker A.
\]

把 $A$ 限制在 $\operatorname{Im}B$ 上，对这个限制映射使用 rank--nullity：
\[
\boxed{
r(AB)
=r(B)-\dim\bigl(\operatorname{Im}B\cap\ker A\bigr).
}
\]

这条式子一次生成所有常用边界。

\subsection{上界}
显然
\[
r(AB)\le r(B).
\]
同时 $AB$ 的 image 位于 $A$ 的 image 中，所以
\[
r(AB)\le r(A).
\]
故
\[
\boxed{r(AB)\le\min\{r(A),r(B)\}.}
\]

\subsection{Sylvester 下界}
中间空间维数为 $n$，而
\[
\dim\ker A=n-r(A).
\]
所以
\[
\dim(\operatorname{Im}B\cap\ker A)
\le n-r(A),
\]
从而
\[
\boxed{r(AB)\ge r(A)+r(B)-n.}
\]

\subsection{什么时候上下界真正取到？}
主公式
\[
r(AB)=r(B)-\dim(\operatorname{Im}B\cap\ker A)
\]
不仅给范围，也直接给取等条件。

若希望第二段作用\kw{完全不再丢方向}，即
\[
r(AB)=r(B),
\]
必须且只需
\[
\boxed{\operatorname{Im}B\cap\ker A=\{0\}.}
\]

若希望 Sylvester 下界取等，交集必须达到它可能的最大维数 $\dim\ker A=n-r(A)$，即
\[
\boxed{\ker A\subseteq\operatorname{Im}B.}
\]
此时 $B$ 已经产生了 $A$ 所有会抹掉的方向，第二段把这部分自由度全部损失掉。

另一侧若已知 $r(A)\le r(B)$，上界 $r(AB)=r(A)$ 等价于 $A$ 在 $\operatorname{Im}B$ 上已经能产生自己的全部 image；可写成
\[
\boxed{\operatorname{Im}B+\ker A=\mathbb R^n.}
\]
这些条件都比“记住何时取等”更稳定，因为它们只是同一个交集公式的不同边界情形。

\subsection{Frobenius 不等式}
复合结构还能生成更一般的三矩阵乘积下界。对于复合 $ABC$，有：
\[
\boxed{r(ABC) \ge r(AB) + r(BC) - r(B).}
\]
Sylvester 不等式其实是令 $B = I_n$ 时的特殊情况。只要盯住“前面保留的图像，有多少继续被后面的核空间吃掉”，所有秩不等式都能从子空间交集的维数中读出。

\subsection{$AB=0$ 为什么能直接变成 rank 约束}
\[
AB=0
\iff
\operatorname{Im}B\subseteq\ker A.
\]
因此
\[
r(B)\le n-r(A),
\]
即
\[
\boxed{r(A)+r(B)\le n.}
\]

\begin{corebox}{不要背 $r(AB)$ 的一堆孤立公式}
先写
\[
\boxed{\operatorname{Im}B\xrightarrow{A}\operatorname{Im}(AB)}
\]
再问：$B$ 已经产生的方向中，有多少正好落入 $\ker A$？所有乘积 rank 结论都围绕这个交集展开。
\end{corebox}

\section{和、转置与 Gram 型矩阵的 rank}

若 $A,B$ 同型，则每个 $(A+B)x$ 都属于
\[
\operatorname{Im}A+\operatorname{Im}B,
\]
因此
\[
\boxed{r(A+B)\le r(A)+r(B).}
\]
这只是上界；两个 image 大量重叠或相互抵消时，实际 rank 可以更小。

由行秩=列秩，
\[
\boxed{r(A^T)=r(A).}
\]

在实数域，
\[
\ker(A^TA)=\ker A.
\]
因为
\[
A^TAx=0
\Longrightarrow
x^TA^TAx=\lVert Ax\rVert^2=0
\Longrightarrow Ax=0,
\]
反向显然成立。两者列数相同，所以 rank--nullity 给出
\[
\boxed{r(A^TA)=r(A).}
\]
同理
\[
\boxed{r(AA^T)=r(A).}
\]

而且 $A^TA$ 不只是“与 $A$ 同 rank”。对任意 $x$，
\[
x^TA^TAx=\lVert Ax\rVert^2\ge0,
\]
所以它天然是半正定矩阵；进一步
\[
\boxed{
A^TA\succ0
\iff
\ker A=\{0\}
\iff
r(A)=n
}
\]
（这里 $A$ 有 $n$ 列）。也就是说，$A^TA$ 正定恰好表示 $A$ 的输入端没有任何方向被压掉。正定性的完整分类仍由 Topic 06 拥有，本册只提供这个由 kernel 直接生成的接口。

这类结论不是“转置以后 rank 神秘保持”，而是 kernel / image 结构没有产生新的自由度损失。

\section{伴随矩阵的 rank：Topic 02 的恒等式怎样在这里继续生长}

Topic 02 已有
\[
A\operatorname{adj}(A)=(\det A)I.
\]
对 $n\times n$ 矩阵（$n\ge2$），rank 分三层：
\[
\boxed{
r(\operatorname{adj}A)=
\begin{cases}
 n,&r(A)=n,\\
 1,&r(A)=n-1,\\
 0,&r(A)\le n-2.
\end{cases}}
\]

\begin{itemize}
  \item 若 $r(A)=n$，则 $\det A\neq0$，且
  \[
  \operatorname{adj}(A)=(\det A)A^{-1},
  \]
  所以满秩；
  \item 若 $r(A)\le n-2$，所有 $(n-1)$ 阶子式都为 $0$，所以每个 cofactor 都为 $0$，即 $\operatorname{adj}(A)=0$；
  \item 若 $r(A)=n-1$，至少一个 $(n-1)$ 阶子式非零，所以 adjugate 非零；又因 $\det A=0$，
  \[
  A\operatorname{adj}(A)=0,
  \]
  adjugate 每一列都落入一维的 $\ker A$，故 rank 至多 $1$。非零再迫使 rank 恰为 $1$。
\end{itemize}

\section{含参数 rank：为什么只需要盯住结构跳变点}

rank 是整数，不会随着参数连续地“慢慢变化”。它改变的时刻，必然有某个原本非零的关键主元或子式变成零。

例如
\[
A(a)=\begin{pmatrix}1&a\\a&1\end{pmatrix}.
\]
其 determinant 为
\[
1-a^2.
\]
所以
\[
r(A)=2\quad(a\neq\pm1),
\]
而当 $a=\pm1$ 时 determinant 归零，但矩阵并非零矩阵，因此
\[
r(A)=1.
\]

这说明含参数 rank 的核心不是“每个参数值重新算一遍”，而是：\kw{找出哪些参数让当前独立结构失效。}

具体消元时，不能未经分类就除以可能为零的参数因子。本册负责解释为什么参数会制造秩跳变；单一参数 rank 问题族的路线优先级由同目录训练 Markdown 维护，跨多个训练专题稳定复用并经证据验证的控制才晋升学科规则。

\section{概念边界：rank 最容易被什么伪装}

\begin{longtable}{L{2.4cm}L{3.0cm}L{3.8cm}L{3.2cm}L{\dimexpr\linewidth-13.0cm\relax}}
\toprule
\textbf{A $\neq$ B} & \textbf{为什么容易混} & \textbf{真正判据} & \textbf{题面信号} & \textbf{混淆后果}\\
\midrule
rank $\neq$ 非零元素个数 & 都像在“计数量” & rank 只计独立方向 & 稀疏/满元素矩阵 & 用元素多少误判自由度\\
rank $\neq$ 非零行数 & 阶梯形里两者恰好相同 & 只有化到阶梯形后非零行才自动无关 & 消元、主元 & 未消元就数行\\
行空间 $\neq$ 列空间 & 行秩=列秩 & 相等的是维数，不是集合 & 非方阵、转置 & 把不同环境空间写成同一集合\\
行变换保 rank $\neq$ 保列空间 & 都涉及左乘可逆矩阵 & $\operatorname{Col}(PA)=P\operatorname{Col}(A)$ & 行初等变换 & 错把输出方向当作不变\\
同 rank $\neq$ 相似 & 两者都是矩阵关系 & 对同型矩阵，同 rank 完全分类“等价”；相似还保谱 & $PAQ$ vs $P^{-1}AP$ & 用 rank 判断算子相似\\
$\det A=0$ $\neq$ rank 已知 & determinant 也测退化 & det 只告诉方阵 rank $<n$ & 奇异方阵 & 把 $n-1$ 与更低 rank 混为一谈\\
$r(AB)$ $\neq r(A)r(B)$ & 标量乘法直觉 & 看 $\operatorname{Im}B\cap\ker A$ & 矩阵乘积、$AB=0$ & 背错误乘法公式\\
$AB=0$ $\neq A=0$ 或 $B=0$ & 数域没有零因子 & 只需 $\operatorname{Im}B\subseteq\ker A$ & 零乘积 & 错过非零矩阵乘积为零\\
\bottomrule
\end{longtable}

\section{看到什么信号时调用本册模型}

\begin{itemize}
  \item 看到“秩、极大无关组、主元、最大非零子式”：统一翻译成“独立信息有几维”；
  \item 看到齐次方程自由变量个数：先回到 $n-r(A)=\dim\ker A$，完整解集转 Topic 04；
  \item 看到两个矩阵做行列初等变换：若目标是分类，问它们是否只在独立重编码输入/输出，从而进入矩阵等价；
  \item 看到 $AB$：先画 $\operatorname{Im}B$ 与 $\ker A$，不要先猜 rank 公式；
  \item 看到 $AB=0$：立即翻译为 $\operatorname{Im}B\subseteq\ker A$；
  \item 看到 $A^TA$、$AA^T$：优先比较 kernel，而不是展开元素；
  \item 看到 adjugate 与 rank：先按 $r(A)=n,n-1,\le n-2$ 分层；
  \item 看到含参数 rank：先定位主元/非零子式何时失效，再分支。
\end{itemize}

\section{一页压缩：rank 是一份“自由度账本”}

\begin{center}
\begin{tikzpicture}[node distance=6mm and 7mm]
\node[nodebox] (input) {输入空间\\$\dim V=n$};
\node[nodebox,right=of input] (split) {选择适配基\\kernel + complement};
\node[nodebox,right=of split] (kernel) {被抹除\\$\dim\ker T=n-r$};
\node[nodebox,below=of kernel] (image) {可见输出\\$\dim\operatorname{Im}T=r$};
\node[nodebox,left=of image] (standard) {矩阵等价标准形\\$J_r$};
\node[nodebox,left=of standard] (product) {复合瓶颈\\$\operatorname{Im}B\cap\ker A$};
\draw[flow] (input) -- (split);
\draw[flow] (split) -- (kernel);
\draw[flow] (split) -- (image);
\draw[flow] (image) -- (standard);
\draw[flow] (standard) -- (product);
\end{tikzpicture}
\end{center}

最终压成四句话：
\[
\boxed{r(A)=\dim\operatorname{Im}A.}
\]
\[
\boxed{n=r(A)+\dim\ker A.}
\]
\[
\boxed{A\sim_{\mathrm{eq}}B\iff r(A)=r(B)\quad\text{(同型矩阵)}.}
\]
\[
\boxed{r(AB)=r(B)-\dim(\operatorname{Im}B\cap\ker A).}
\]

如果忘掉具体公式，重新问：\kw{输入有多少自由度？哪些方向被压没？哪些方向还能被输出看见？复合时中间 image 有多少撞进下一步 kernel？}

\end{document}
